\documentclass[10pt,twocolumn]{article}
\usepackage[margin=0.78in]{geometry}
\usepackage[T1]{fontenc}
\usepackage{lmodern,microtype,amsmath,amssymb,amsthm,booktabs,graphicx}
\usepackage[raggedright]{titlesec}
\usepackage[round,authoryear]{natbib}
\usepackage{xcolor,xurl}
\usepackage[colorlinks=true,allcolors=blue]{hyperref}
\usepackage{caption}
\captionsetup{skip=6pt}
\captionsetup[table]{position=top}
\setlength{\columnsep}{0.25in}
\setlength{\parindent}{1em}
\setlength{\parskip}{0pt}
\setlength{\emergencystretch}{1.5em}
\setlength{\bibsep}{0pt plus 0.3ex}
\ifdefined\layoutfixture
\usepackage{fancyhdr}
\pagestyle{fancy}\fancyhf{}
\fancyfoot[C]{\scriptsize LAYOUT CHECK: SYNTHETIC NUMBERS \quad \thepage}
\renewcommand{\headrulewidth}{0pt}
\fi
\newtheorem{proposition}{Proposition}
\newcommand{\E}{\mathbb{E}}
\newcommand{\R}{\mathbb{R}}
\input{\paperroot/metrics.tex}
\input{\paperroot/tables.tex}
\title{Auditing Reference Distributions for\\ Observation-Replacement Return Spans}
\author{Aaryan Singh\\
\small\texttt{singha9@rose-hulman.edu}\\
\small Code and data: \url{https://github.com/Aaryansi/nano-l1}}
\date{}
\begin{document}
\raggedbottom
\maketitle
\ifdefined\layoutfixture
\begin{center}\fbox{\textbf{LAYOUT CHECK: SYNTHETIC NUMBERS}}\end{center}
\fi
\begin{abstract}
An observation-replacement return span measures how a specified input
intervention changes a policy's return. Interpreting that span requires choices
about the reference population, replacement distribution and evaluation process.
We audit these choices in a binary-contract simulator, controlled synthetic
populations and six control-task configurations. We retain training-seed
variation, raw endpoint returns, failed competence gates and unresolved
references. In the historical market example, five policies have mean return
\MarketMean\ dollars per episode with a crossed seed/day interval
\MarketInterval. Reference comparisons remain descriptive because market
policies are not exchangeable with the synthetic bank. Across \BankCount{}
independent synthetic reference banks, the zero-strength flag rate is
\BankNullRate{} for each rule, but its conservative bank-level upper bound is
\BankNullUpper. This does not certify five-percent error control.
The zero-strength population still has positive full-input replacement spans,
illustrating that return sensitivity can arise without settlement-predictive inputs.
An exact finite-pool example shows why the legacy reference-centered
rank/normal rule lacks a general exchangeability guarantee; standard symmetric
ranks address that separate statistical requirement. The audit supports
careful interpretation of particular interventions, without inferring
feature-level fidelity, absent information or explanation meaning from a
single return-span comparison.
\end{abstract}

\section{Question and scope}
An explanation can describe a policy's decisions, its predictions, or the
consequences of its actions. These are different targets, including when the
policy performs poorly. SVERL makes this distinction explicit and assigns
feature contributions to specified characteristic functions
\citep{beechey2023,beechey2025}. A faithful explanation of a costly decision can
be useful for debugging it. Low return alone does not make the explanation
empty.

Our question concerns a narrower diagnostic: how should an observed
observation-replacement return span be compared with spans from a bank of
reference policies? We examine the implemented comparison, the financial state
retained by an intervention, and the evidence needed to interpret reference
widths and classifications. The contributions are a reproducible computational
audit, explicit examples of interpretation failures, and empirical comparisons
under declared training budgets. We do not propose a universal null for
reinforcement learning or a new general rank-test theorem.

The study is retrospective. Its configuration was fixed for the selected
rerun after exploratory development on this historical data. Repeating that
window does not create a prospective holdout. All numerical tables are
generated from the selected primary release and two separately identified
supplemental studies; the primary training source is
\texttt{\TrainingCommit}. Reduced-budget execution runs cannot supply these
tables.

\section{What the endpoint statistic identifies}
Let $\pi$ be a fixed policy, $M$ an environment, $D$ an evaluation population,
and $I$ a specified replacement intervention. For a feature set $S$, let
$v_{\pi,M,D,I}(S)$ denote expected episode return when the intervention retains
the coordinates in $S$. The endpoint span is
\begin{equation}
\Delta(\pi,M,D,I)=v(N)-v(\emptyset).
\label{eq:span}
\end{equation}
We estimate both endpoints by replaying the same held-out episodes. The full
endpoint supplies original policy inputs; the empty endpoint applies the
declared replacement at every decision. The intervention also changes the
policy's actions and subsequent inventory. Therefore this is a property of the
policy, environment, intervention and evaluation distribution together.

For an efficient Shapley decomposition, $\sum_i\phi_i=v(N)-v(\emptyset)$
\citep{lundberg2017}. Reading the endpoints avoids estimating individual
coordinates merely to recover their sum. It does not validate that
decomposition: allocations with the same sum receive the same diagnostic
regardless of how they assign feature credit. Estimation accuracy for a
ranking is another question; our implementation's gap diagnostic does not
implement the formal ranking procedures of \citet{goldwasser2024}.

The intervention is also distinct from SVERL's conditional policy-occupancy
construction. The extended framework defines outcome explanations through
partial-information behaviour at the state being explained, with the original
policy elsewhere \citep{beechey2025}. Our empirical flat-inventory background
and repeated replacement at all steps differ in distribution and locality.
The experiments do not test or refute that theoretical construction.
FastSVERL provides practical approximations for these distinct behaviour,
outcome and prediction targets, including sequential and off-policy estimation
\citep{beecheyfast2025}. Our endpoint audit does not benchmark that method.
Background selection and the relationship between SHAP feature importance and
agent performance have also been studied empirically, including exceptions in
fragile tasks \citep{engelhardt2024}. We claim neither a new discovery of
background sensitivity nor a general equivalence between importance and return.

\ifdefined\shortversion
Parameter sensitivity and probability attribution address separate targets;
completeness does not make a probability difference a return intervention
\citep{adebayo2018,sundararajan2017}.
\else
\subsection{Dependence, fidelity and task utility}
Parameter-randomization sensitivity asks whether an explanation depends on
learned parameters \citep{adebayo2018}. It does not, by itself, establish
feature-level fidelity or useful task performance. Limitations of such checks
and refinements to their evaluation are established topics
\citep{binder2022,hedstrom2024}. A bank of independently initialized policy
spans is a different experiment from progressively reinitializing the layers
of one trained model.

Integrated gradients satisfy completeness for their chosen model output
\citep{sundararajan2017}. Our auxiliary IG implementation explains a
chosen-action probability relative to a zero normalized input. That
probability difference is not the return intervention in
Equation~\ref{eq:span}. We therefore do not claim cross-family return-test
agreement from these auxiliary profiles.
\fi

\section{Reference diagnostics and rank validity}
The legacy rule compares an observed span $T$ with $B$ reference spans
$T_1,\ldots,T_B$. It combines a reference-centered empirical tail score
\begin{equation}
p_{\rm legacy}=\frac{1+\sum_{b=1}^B
\mathbf{1}\{|T_b-\bar T_{\rm ref}|\geq|T-\bar T_{\rm ref}|\}}{B+1}
\label{eq:legacy}
\end{equation}
with a two-sided normal approximation to
$z=(T-\bar T_{\rm ref})/s_{\rm ref}$. Both criteria must meet the declared
level $\alpha=0.05$. We retain this rule as a diagnostic of the original
workflow, not as an assumption-free test. Insufficient rank resolution or
$s_{\rm ref}\leq10^{-12}$ makes its comparison unresolved. This convention
does not mean that every test against a point-mass reference is invalid.

\begin{proposition}
The legacy combined rule does not have a general finite-sample
exchangeability guarantee.
\end{proposition}
\begin{proof}
Uniformly designate one observation in a fixed pool containing $-1,+1$ and
18 zeros; the remaining 19 entries are references. The resulting observation
and references are jointly exchangeable. Either extreme receives legacy tail
score $1/20=0.05$ and $|z|=20/\sqrt{19}>4.58$, so both meet the normal
criterion. The rejection probability is $2/20=0.10$. This is a finite-pool
example, not an error estimate for the independent synthetic benchmark.
\end{proof}

For comparison, let $T_0=T$ and compute the pooled center
$\bar T_{\rm pool}=(B+1)^{-1}\sum_{b=0}^{B}T_b$. The symmetric two-sided
rank uses scores $A_b=|T_b-\bar T_{\rm pool}|$ and
\begin{equation}
p_{\rm sym}=\frac{\sum_{b=0}^{B}\mathbf{1}\{A_b\geq A_0\}}{B+1}.
\label{eq:symmetric}
\end{equation}
We also report the upper-tail raw-span rank. Both include the observation and
count ties conservatively. Under joint exchangeability and permutation-equivariant
scoring, their ranks are super-uniform by the established finite-rank argument
\citep{kuchibhotla2020,hemerik2018}. Conditional on an unordered score multiset,
at most an $\alpha$ fraction of indices can have an inclusive upper-tail rank
at most $\alpha$; exchangeability makes the observed index uniform.
The symmetric rank assigns each extreme in the example $2/20=0.10$.

That guarantee is marginal over the bank and the observed draw. It does not
bound error conditional on each realized reference bank. It also does not
supply exchangeability between a historical market agent and policies trained
on synthetic populations or corrupted inputs. Connections between attribution
and conditional-independence testing require explicit modeling assumptions,
as in SHAP-XRT \citep{teneggi2023}; the return setting does not inherit those
guarantees automatically.

\section{Implemented experiments}
\subsection{Market accounting and evaluation}
The pinned corpus contains \EpisodeCount\ binary-contract episodes with 14
decision rows each. Chronological training, validation and test counts are
\TrainCount, \ValidationCount\ and \TestCount, with four boundary episodes
purged. The 18 inputs comprise 14 market/spot coordinates normalized using
training statistics and four active-inventory coordinates. Actions target
$-100,0,+100$ signed yes-equivalent contracts. A reversal can trade 200.

The full-input replacement background contains 256 rows sampled from all
evaluation steps while inventory is held flat. Each intervened episode draws
a fresh whole background row at every decision. This replaces the clock as
well as market and inventory coordinates; it is an empirical intervention,
without a guarantee of compatibility with the recipient's state.

The assumed execution model buys at the supplied ask and sells at the bid.
It charges the per-order fee
$\operatorname{ceil}_{\rm cents}[0.07|C|p(1-p)]$, with zero impact by default,
full fills and no funding or queue simulation. Candle-derived quotes are
coarse proxies; these assumptions do not establish exact historical fees or
live executable performance. The default terminal mark is the recorded binary
settlement, with no additional settlement fee in the model. Equity increments
telescope to final profit and loss at $\gamma=1$.

Five market seeds train for 100 PPO updates under each cost objective. Their
greedy policies are evaluated on every held-out episode. The critic predicts
stochastic policy return in dollars, not settlement probability; its auxiliary
calibration uses eight policy rollouts per episode. Both rollout implementations
agree on every acted-on observation and greedy action for all ten checkpoints,
with maximum accounting discrepancy \ParityError. Primary uncertainty crosses
training-seed resampling with shared UTC-day blocks. There are only
\DayCount\ test-day blocks, limiting inference about temporal variation.

\subsection{Synthetic populations and reference banks}
The balanced synthetic population has independent binary settlements, constant
price 0.5, spread 0.02, and noisy features independent of settlement. Flat earns
zero. The planted population adds a noisy settlement signal, repeated across
the 14 steps. Its first-observation sign-and-hold strategy is a reference
policy, not a proven optimum when aggregation, delayed entry and abstention
are possible.

The primary bank has 24 independently seeded policies, each trained for 40
updates on 1,200 synthetic episodes and evaluated on an independent draw of
1,200 episodes with its own training normalizer. A 99-draw sensitivity retains
the same procedure. The conditional calibration uses 30 independent observed
trials at each strength $0,0.2,1$, reusing the primary bank. Each trial retains
both endpoint returns. We report pointwise Wilson intervals conditional on
that bank, alongside unresolved counts and both standard rank comparators.
These rates include the specified learning procedure and budget. They do not
isolate test sensitivity from learning success or quantify bank-to-bank error.

A separate study repeats the experiment across \BankCount{} independent banks,
each with 24 references and 30 observed trials at each of the three strengths.
All policies train for 40 updates on fresh 1,200-episode corpora, use their own
training normalizer and evaluate on a fresh 1,200-episode corpus. Every policy
has distinct data, training and masking seeds. Unlike the original conditional
study, masking seeds are separate from policy seeds. This is a declared common
procedure for all three rules, not a comparison isolating bank reuse alone.

The independent units are $K=\BankCount$ banks. For bank flag fractions
$X_k\in[0,1]$, we estimate the marginal flag probability by their equally weighted mean. We
invert the bounded-variable Chernoff bound \citep{hoeffding1963} with threshold
$\log(2/0.05)/K$ for Bernoulli KL divergence. This yields conservative pointwise
95\% intervals over banks, including at zero or unit observed rates. It does
not treat the 600 trials as independent Bernoulli observations. Paired method
differences use the same bound after mapping $[-1,1]$ to $[0,1]$. Intervals are
not simultaneous over rules, strengths or contrasts. Unresolved comparisons
remain in the operational flag-rate denominator and are reported separately.

\subsection{Control tasks and competence gates}
We evaluate CartPole, Acrobot, MountainCar, Pendulum with nine discrete torques,
deterministic $4\times4$ FrozenLake, and CliffWalking capped at 200 steps.
The control trainer is Stable-Baselines3 PPO \citep{raffin2021}, with separate
two-layer 64-unit actor/value trunks and four environment copies. Observed
agents request 400,000 interactions; each blinded reference requests 30,000.
Complete rollout batches can exceed those requests. Each construction has 24
reference draws, and endpoint evaluation uses 20 episodes. Five independently
trained observed seeds are separate from correlated training checkpoints.

The parameter reference uses fresh orthogonally initialized actors, with no
learning. The environment reference trains agents with independent
coordinate-wise Gaussian observations whose moments are estimated from a
random-policy rollout, leaving transitions and rewards intact. Both are
evaluated on the original task with full inputs and with repeated background
replacement; the background contains 512 random-policy observation rows.
Gaussian blinding can violate within-state constraints. A separate whole-row
resampling sensitivity preserves donor row support, without guaranteeing
consistency with the recipient trajectory. Initialization and blinded learning
are different reference procedures, without an established exchangeability
relationship to competent observed agents.

Declared competence thresholds are respectively 475, $-100$, $-110$, $-300$,
0.8 and $-20$, measured on 25 episodes, together with an absolute measured
endpoint span above $10^{-6}$. A failed gate does not establish low detection power. Neither
reference-width ratios nor standardized distances are power multipliers.
All failed and unresolved cases remain in the results.

\section{Empirical results}
\subsection{Historical market and primary references}
\begin{table}[t]
\centering\small
\caption{Held-out market return in dollars per episode. Seed spread is the
sample SD across seed means, not episode SD or a standard error. Cost and
zero-cost agents were trained separately. The primary crossed seed/day
interval is \MarketInterval. Non-rejection would not establish equivalence.}
\label{tab:market}
\MarketTable
\end{table}
The market policies have mean return \MarketMean\ dollars per episode
(Table~\ref{tab:market}). The primary centered-bootstrap probability is
\MarketProbability. This approximate probability and the percentile interval
are different summaries and need not be dual near the decision level. We
report both, without selecting one to support a stronger claim. Neither a
negative mean nor a comparison with flat proves absence of predictive
information or makes a policy's explanation meaningless.

\begin{table}[t]
\centering\small
\caption{Primary synthetic-bank diagnostics. ``Outside'' denotes the legacy
two-part classification, not a finding of explanation fidelity. Market
comparisons are descriptive. Undefined distances are shown explicitly.}
\label{tab:headline}
\HeadlineTable
\end{table}
Table~\ref{tab:headline} compares the fixed primary cases against the declared
bank. All five market-seed comparisons and the larger-bank sensitivity are
retained in the release. \MarketSeedParagraph{}
\SensitivityParagraph{}
Changing a bank can change its center, width and
classification while also changing training or evaluation distributions.
Those differences do not identify a unique information-removal mechanism.

\subsection{Conditional and marginal calibration}
\begin{figure}[t]
\centering
\includegraphics[width=\columnwidth]{\paperroot/figures/calibration.png}
\caption{Conditional trial rates with pointwise intervals for one reused
reference bank. A standard rank guarantee under exchangeability is marginal
over banks; these intervals do not establish such a guarantee on the market.}
\label{fig:calibration}
\end{figure}
\CalibrationParagraph{} Figure~\ref{fig:calibration} plots these rates and
their intervals.
With 30 trials, even zero flags leave a substantial upper uncertainty bound.
One positive-control success is therefore not a power estimate. The signal
strength sweep also changes what the learning procedure can acquire within
its fixed budget.

\begin{table*}[t]
\centering\small
\caption{Independent-bank calibration: \BankCount{} banks with 30 trials per
strength in each bank. Rates and pointwise bounded-bank intervals are rounded
to three decimals. All declared trials remain in the denominator. Unres. counts
unresolved comparisons. The signal alternatives measure learning and comparison
together, not test power in isolation.}
\label{tab:banks}
\BankTable
\end{table*}
Table~\ref{tab:banks} adds bank-to-bank variation to the audit. The zero-strength
rate is \BankNullRate{}, but its upper bound is \BankNullUpper. The observed
rate below 5\% does not certify error control at 5\%. \BankAgreementParagraph{}
This empirical agreement does not contradict the finite-pool counterexample:
that example rules out a general guarantee for the legacy rule, rather than
predicting its error in every population. The bank study retains all \BankModelCount{}
checkpoints and update logs; both endpoints were reproduced from every saved
policy after extraction from the distributed archive.

\subsection{Matched objective replication}
The same market episodes and state transitions support two training objectives:
trading profit in dollars, or the sum of signed settlement calls at each step,
with no execution costs in the prediction score. Both objectives train from
scratch at 40 and 100 updates. Each cell has five observed policies and 24
Gaussian-blinded references trained under that cell's objective and budget.
The normalizer uses the original training split, and every policy is evaluated
on the complete test partition with both full and replaced observations.
This study uses a 192-row flat-inventory background, common masking draws
across observed agents and reference-specific masking draws. Its replacement
background and draws differ from the primary 256-row comparison, so those
span differences cannot be attributed to the objective alone.
\begin{table}[t]
\centering\small
\caption{Matched-budget objective cells. Return and span are means over five
observed seeds. Reference SD describes 24 blinded policies. Prediction uses
signed-call score units; trading uses dollars. Cross-objective magnitudes and
standardized distances do not measure a causal objective effect.}
\label{tab:objectives}
\ObjectiveTable
\end{table}
\ObjectiveDecisionParagraph{} Table~\ref{tab:objectives} shows why that shared classification does
not identify a shared mechanism: reference widths and observed spans change
with the task and budget. The supplemental replication retains raw full and
masked endpoints for \ObjectiveObservedCount{} observed and
\ObjectiveReferenceCount{} reference policies. \ObjectiveMatchParagraph{} This reconstructs
the new replication's masked endpoints; it does not recover missing original
arrays or replay original objective checkpoints. No objective-effect
significance or explanation-fidelity claim follows from these comparisons.

\subsection{Control-task competence and references}
\begin{table}[t]
\centering\footnotesize
\caption{Independent control seeds. Eligibility requires both the declared
return gate and measured endpoint dependence. Classification counts retain
all seeds and report unresolved references separately. Each task uses its
native reward units. O/U denotes outside/unresolved counts out of five seeds.}
\label{tab:controls}
\ControlTable
\end{table}
\ControlParagraph{}
Table~\ref{tab:controls} supports statements about these configurations and
reference banks.
It does not establish a universal ranking of reference constructions. A
reference can collapse when its policies have identical returns, and the
legacy normal-scale rule then cannot decide. That outcome differs from a
resolved non-rejection and from a trained observed agent failing to learn.

\subsection{State preservation and attribution suppression}
The market state controls use the five saved cost-enabled checkpoints. Each
synthetic case trains five new policies for 60 updates, using 1,200 episodes
per training draw and 1,200 common independent evaluation episodes.
Full replacement changes inventory inputs as well as market inputs. The
market-row and donor-trajectory controls retain the real clock and all four
inventory coordinates of the intervened rollout. A donor row preserves empirical row support, and a
donor trajectory preserves its own market sequence. Neither guarantees joint
consistency with the recipient's financial state or reward.

\begin{figure}[t]
\centering
\includegraphics[width=\columnwidth]{\paperroot/figures/state_interventions.png}
\caption{Market return spans under three specified interventions. Intervals
cross training seeds and common UTC-day blocks, conditional on fixed masking
draws. Differences do not isolate inventory as their sole cause.}
\label{fig:state}
\end{figure}
\StateParagraph{} Figure~\ref{fig:state} compares the three interventions.
These runs fix masking draws within each analysis; their recorded seeds differ
from the primary comparison's draws for some policies. \NullStateParagraph{}
Positive sensitivity in this generator need not indicate settlement-predictive
information: replacement can change trading and its costs. These controls do
not isolate one cause of the span differences.

Retaining inventory also retains some price information. If $z=q/100$ is
scaled nonzero position, $e$ its entry price and $u=q(P-e)/100$ scaled
unrealized profit, then
\begin{equation}
P=e+u/z.
\end{equation}
The identity holds up to feature precision. Thus a financially consistent control retaining these coordinates cannot
also hide all current-price information when $z\neq0$. Conditional information
beyond retained inventory can still be a legitimate target; this identity is
not an impossibility result for conditional explanations.

The steering experiment penalizes input gradients during training and measures
the resulting behaviour-attribution share and held-out return separately.
An auxiliary single-seed comparison examines greedy behaviour; it cannot
establish behavioural equivalence across the full experiment.
The penalty targets the expiry clock in the market and the planted
signal in the synthetic case, with ten seeds and penalty 20 fixed for the
reported comparison. Each steering policy trains for 60 updates. Synthetic
training and independent evaluation each use 1,500 episodes; market evaluation
uses the complete test partition. All penalties are retained and shown in
Figure~\ref{fig:steering}. \SteeringParagraph{}
\begin{figure}[t]
\centering
\includegraphics[width=\columnwidth]{\paperroot/figures/steering.png}
\caption{All declared penalties. Upper panel shows seed-specific target shares
and their means; lower panel shows paired return changes with crossed seed/block
intervals. The two populations target different features.}
\label{fig:steering}
\end{figure}
A non-significant return contrast is not equivalence or noninferiority, and
changed task return alone cannot establish that the resulting policies are
behaviourally identical.

\ifdefined\shortversion\else
\section{Why settlement shuffling changes the question}
Global settlement permutation preserves feature arrays and the marginal
settlement rate. It changes their relationship to rewards. This matters even
in a calibrated one-period example with no incremental price-measurable edge.
Suppose $Y\in\{0,1\}$, $P\in[0,1]$ and $\E[Y\mid P]=P$. Every bounded
price-measurable position $q(P)$ has frictionless gross expected payoff
\begin{equation}
\E[q(P)(Y-P)]=0.
\end{equation}
If $Y^*$ is replaced by an independent draw from its marginal settlement law,
with mean $\mu=\E[P]$, choosing
$q^*(P)=\operatorname{sign}(\mu-P)$ instead gives
\begin{equation}
\E[q^*(P)(Y^*-P)]=\E|\mu-P|,
\end{equation}
which is positive for a nonconstant price. Fees and execution constraints can
reduce this gross payoff. The example does not establish live profit or
efficiency of the historical corpus. It shows that independence of settlement
and price is a different hypothesis from calibrated prices with no additional
edge. Price remains reward-relevant through acquisition cost even after losing
settlement-predictive information.

Within-price-bucket label changes and correlations are descriptive diagnostics,
not estimates of conditional information removal. A conditional null may
legitimately retain predictive information in its conditioning price. Failure
to meet a finite set of heuristic bucket thresholds does not prove conditional
randomization impossible or vacuous. A sequential test would need a justified
joint conditional law for paths and outcomes.
\fi

\section{Limitations and next experiments}
The financial model uses assumed costs and coarse quotes, the market period
is fixed, and seed/day inference rests on few time blocks. Synthetic
calibration includes a conditional study and an independent-bank study, but
still has limited precision at rare error rates. Control competence is
configuration-specific; failed tasks remain
evidence about training rather than about detection sensitivity. No reference
width alone identifies an explanatory mechanism.

\ifdefined\shortversion
Complete replay and hash checks verify computations, without establishing
feature-level fidelity. Additional raw evidence and implementation limits
remain documented in the full profile and release.
\else
This audit distinguishes computational correctness from scientific
identification. Complete replay parity, corrected KernelSHAP coalition-size
mass and exact-game tests strengthen the implementation evidence. They do
not turn the endpoint sum into a feature-fidelity test. Auxiliary ranking,
IG and parameter-randomization diagnostics address different targets and
assumptions. The same applies to secondary reference sensitivities: some masked
spans retain means without per-episode masked arrays. The objective supplement
now retains both endpoints and reproduces the original summaries, but its
policies have no archived checkpoints or training logs. Historical evidence
retains its original scope rather than acquiring the replication's additional
evidence retroactively.
\fi

The next useful experiments are competence-controlled reference comparisons
with matched learning budgets and conditional path randomization in a simulator
with a known joint law. Such a benchmark should assess individual feature
allocations, including deliberately incorrect allocations with the same
endpoint sum. A larger independent-bank study should choose its sample size
and precision target before further runs; extra trials within existing banks
do not create additional independent banks.
A separately frozen temporal market holdout would test external
robustness. These additions would strengthen calibration and generalization;
they are not prerequisites for honestly reporting the present descriptive
audit.

\paragraph{Reproducibility.}
The repository pins the corpus and splits, retains raw returns and model
hashes, records numerical runtime and package versions, and verifies source
against the training commit. The full runner and release verifier are separate
from execution-only smoke checks. Manuscript tables, figures and summaries
share one derived statistics artifact; the three document profiles share this
source.

\ifdefined\extendedversion
\appendix
\section{Additional numerical records}
\begin{table}[!ht]
\centering\small
\caption{Every market-seed comparison against the primary bank. These remain
descriptive comparisons with a synthetic reference population.}
\MarketSeedTable
\end{table}
\begin{table}[!ht]
\centering\small
\caption{Descriptive cost-enabled baseline returns on the same held-out
episodes. A lower mean is not a statement about which model has a faithful
explanation. The momentum fallback is labeled explicitly.}
\BaselineTable
\end{table}
\begin{table}[!ht]
\centering\small
\caption{All steering penalties. Attribution shares concern the specified
behaviour proxy. Return differences use common evaluation blocks and paired
training seeds. An interval containing zero does not establish equivalence.}
\SteeringTable
\end{table}
\begin{table}[!ht]
\centering\small
\caption{State-preservation spans and crossed seed/block intervals.}
\StateTable
\end{table}
\paragraph{Stored auxiliary analyses.}
The release also retains the reference-training budget sweep, the larger-bank
sensitivity, nearest-neighbor masking, explanatory target proxies, layer
randomization and horizon-width studies. Nearest-neighbor replacements do not
provide a known conditional law. Horizon lengths change clock inputs and
trading opportunities together; their log-log fits are descriptive. Unresolved
comparisons and undefined correlations are retained rather than converted to
agreement, no signal or equivalence.
\fi
\bibliographystyle{plainnat}
\bibliography{\paperroot/references}
\end{document}
